\section{Conclusão}

\textbf{O que foi encontrado.} No simulador construído com dados de uma empresa de transporte por aplicativo que opera na Grande Vitória (ES), distribuir em conjunto os pedidos de cada lote de 60 segundos, pelo algoritmo Húngaro, gastou menos km vazio que mandar o veículo mais próximo assim que o pedido chega: 5,96\% a menos em setembro e 5,04\% a menos em outubro de 2026, com a configuração congelada antes. Isso equivale a cerca de 1,3 e 0,9 ponto percentual do km total modelado. A redução apareceu em todas as variantes de sensibilidade e em todos os meses de fevereiro a agosto. Nos cenários sintéticos, exploratórios, a redução mediana foi de 2,0\% a 8,9\%.

\textbf{O que não foi confirmado.} Em outubro, a atribuição conjunta atendeu 85 pedidos a menos que o despacho imediato e não cumpriu o critério de atendimento; ela atendeu menos em 10 das 12 variantes desse mês. O critério cumprido em setembro vale com os parâmetros de março: com a calibração refeita, ela atendeu 40 pedidos a menos também em setembro. Contra o lote simples, que espera o mesmo lote mas decide um pedido por vez, a vantagem foi de 2,74\% e 1,40\% nos períodos congelados, mas pequena e frágil: de −0,3\% a −4,6\% nas variantes. No único teste fora da amostra, 12 dias de 26/10 a 08/11 com calibração refeita, a redução contra o despacho imediato caiu para 2,79\%, sem se distinguir de zero, e o km ficou 0,87\% acima do lote simples, embora a atribuição conjunta tenha atendido 29 pedidos a mais, com o mesmo km por atendido (2,84). O teste tem pouco poder, mas não confirma a vantagem sobre o lote simples. Um guloso global simples, definido depois, teve menos km que a atribuição conjunta e atendeu bem menos com os parâmetros de março (179 e 83 pedidos a menos); com a calibração refeita, essa diferença caiu para 41 e 25 e, no teste, se inverteu (29 atendimentos a mais e 0,62\% menos km). A vantagem do Húngaro sobre esse guloso não se confirmou. O simulador também não reproduz a operação: atende de 9 a 23 pontos percentuais a mais, mesmo depois de recalibrado. Assim, o estudo compara políticas no modelo e não demonstra economia na operação real nem generaliza para outras empresas ou cidades.

\textbf{O que faltaria para saber se funciona na operação.} Modelar o aceite e a recusa de corridas pelos motoristas; validar o simulador contra a operação observada; comparar com o despacho que a empresa de fato pratica, com distâncias pela rede viária; e, por fim, fazer um piloto controlado na operação.

\section*{Declarações}

\textit{Dados.} As entradas históricas vêm de uma empresa de transporte por aplicativo que opera na Grande Vitória (ES). Foram extraídas em modo somente leitura e processadas localmente. Identificadores de motoristas, veículos e passageiros foram pseudonimizados por HMAC-SHA256 com segredo mantido fora do repositório, e só agregados são reportados.

\textit{Código e dados.} O código, as configurações e os resultados agregados estão no repositório dos autores, não publicado. Os dados individuais não podem ser distribuídos. Terceiros podem reproduzir os cenários sintéticos (geografia pública), mas não a avaliação com dados reais sem acesso aos dados.

\textit{Registro.} O congelamento foi documentado por \textit{hashes} locais, sem registro externo com carimbo de tempo; sua data não é verificável por terceiros. As análises da Subseção~\ref{sec:validade}, inclusive a recalibração e o teste de 26/10 a 08/11, e a dispersão da Tabela~\ref{tab:sint} foram feitas depois dos resultados.
